\begin{conclusion}

本文围绕移动机器人在复杂环境下的自主导航问题，完成了理论梳理、多传感器融合定位与
路径规划避障三部分工作。理论方面汇总了神经网络、卷积与循环网络以及强化学习的数学
表述；定位方面建立了多传感器的时间同步与外参标定流程，采用误差状态卡尔曼滤波融合
激光雷达、惯性测量单元与轮式里程计，实验表明定位误差较单一传感器降低三成以上；
规划方面在 A$^{*}$ 算法的代价函数中引入转向惩罚项，并在动态窗口法中加入全局路径
引导项，仿真显示狭窄通道下的通行成功率明显提高。相对于任务书的要求，本文的实验
均在室内结构化环境完成，未引入视觉信息，动态障碍的运动规律也与真实行人存在差距。
后续可引入视觉惯性里程计、把方法扩展到多机器人协同，并研究外参在线标定与参数选取
的理论依据。

\textcolor{red}{说明：总结与展望，总结课题完成了哪些相关工作，得到的成果如何（或得到了什么结论），相对于任务书要求，还存在哪些不足；可以提出建议、研究设想、改进思路、尚待解决的问题等等。}

\end{conclusion}
